% ===================================================================
%  ТАБЛИЦЯ № 5 — Дім і як його будують.
%  Traduction 2026 d'apres texto/io, controlee sur texto/fr, texto/en
%  et texto/es. Consigne : voir texto/uk/10-tablycia-01.tex.
%
%  Le tableau 5 est un DIALOGUE, et les deux volumes composent le nom
%  du parleur en petites capitales : on garde \textsc{}, que le rendu
%  laisse tomber en gardant le texte.
%
%  Les renvois du plan sont des LETTRES — (a) le couloir, (b) le
%  salon... — et non des numeros. On les ecrit comme l'ido les ecrit,
%  y compris la ou l'imprimeur a compose « (1) » pour « (l) ».
%
%  LE TABLEAU DES METIERS. Les noms de corps de metier sont la ou les
%  deux langues slaves divergent le plus franchement : le macon est un
%  « муляр » et non un « каменщик », le charpentier un « тесля » et non
%  un « плотник », le platrier un « тинькар » et non un « штукатур »,
%  le zingueur un « бляхар » et non un « жестянщик ». Quatre metiers
%  du meme chantier, quatre mots sans parent commun.
%
%  LES NOMS D'HOMME RESTENT FRANCAIS. La colonne russe a russifie
%  l'entrepreneur — « Blanc » y devient « Белов » — et c'est un parti
%  qu'on ne reprend pas : on translittere, « пан Блан », comme on
%  translittere « пан Ру » et « пан Бернар ».
% ===================================================================

%%K t05-apar-1 apar
\VUsaut{6.0mm}
\VUtitre{15.0pt}{60}{ТАБЛИЦЯ № 5}
\VUsaut{1.4mm}
\VUtitre{11.4pt}{0}{Дім і як його будують.}
\VUsaut{2.2mm}

%%K t05-tit-1 sub
\VUpk{10.2pt}{40}{Щаслива зустріч.}
\VUsaut{3.0mm}

%%K t05-01-1 p
\textsc{Іван}. --- Дядьку, мені дуже хотілося б дізнатися, як будують
\VUgras{дім}.

%%K t05-01-2 p
\textsc{Дядько} \textsuperscript{(80)}. --- Це неважко, мій хлопче; ходімо зі
мною, і я покажу тобі той, що будує пан Бернар \textsuperscript{(79)} і в якому я
житиму.

%%K t05-01-3 p
\textsc{Іван}. --- А хто накреслив \VUgras{план} \textsuperscript{(76)} цього
дому?

%%K t05-01-4 p
\textsc{Дядько}. --- \VUgras{Архітектор} \textsuperscript{(75)}; він зробив дуже
ясне креслення, його схвалили, і \VUgras{підрядник} \textsuperscript{(77)}
розпочав роботи.

%%K t05-01-5 p
\textsc{Іван}. --- А хто підрядник?

%%K t05-01-6 p
\textsc{Дядько}. --- Пан Блан, батько твого товариша Якова. Він керує
робітниками разом із паном Ру, архітектором.

%%K t05-01-7 p
Та ось і він! Пан Блан іде якраз вчасно зі своєю \VUgras{лінійкою}
\textsuperscript{(78)}, а пан Ру зі своїм планом.

%%K t05-01-8 p
Доброго дня, панове. Як ся маєте?

%%K t05-01-9 p
\textsc{Пан Ру}. --- Чудово, дякую. Добридень, Іване; як хлопець виріс!

%%K t05-01-10 p
\textsc{Дядько}. --- І справді, зовсім маленький чоловік. Він дуже поважний;
йому треба пояснювати все, що він бачить. Сьогодні йому хочеться дізнатися, як
будують дім.

%%K t05-tit-2 sub
\VUpk{10.2pt}{40}{Робітники.}
\VUsaut{3.0mm}

%%K t05-01-11 p
\textsc{Пан Блан}. --- То ходімо зі мною, мій юний друже, і я тобі зараз же все
поясню, бо дуже люблю дітей, які хочуть учитися.

%%K t05-01-12 p
Бачиш того робітника з \VUgras{лопатою} \textsuperscript{(82)} в руках: це
\VUgras{землекоп} \textsuperscript{(81)}. Він риє \VUgras{канаву}
\textsuperscript{(46)}, щоб покласти водогінну трубу. Він же вирив землю там, де
стоїть дім, щоб муляр міг вивести чотири \VUgras{стіни}. Та частина цих стін, що
в землі, зветься \VUgras{фундаментом} \textsuperscript{(2)}. Простір, замкнений
фундаментом, утворює \VUgras{підвал} \textsuperscript{(3)}. У цього підвалу є
віконце, зване \VUgras{продухом} \textsuperscript{(5)}, \VUgras{двері}
\textsuperscript{(4)} і сходи.

%%K t05-01-13 p
\VUgras{Муляр} \textsuperscript{(108)} далі виводив стіни, лишаючи отвори для
\VUgras{дверей} \textsuperscript{(8)} і \VUgras{вікон} \textsuperscript{(17)}.

%%K t05-01-14 p
\textsc{Іван}. --- Але я його не бачу, цього муляра!

%%K t05-01-15 p
\textsc{Пан Блан}. --- Він уже скінчив працювати на домі. Він он там, коло
\VUgras{стайні} \textsuperscript{(54)}, на своєму \VUgras{риштуванні}
\textsuperscript{(110)}, кладе стіну \VUgras{пральні} \textsuperscript{(56)}.

%%K t05-01-16 p
У нього є кельма, ночви і \VUgras{висок} \textsuperscript{(109)}. Але цих назв ти
не запам'ятаєш.

%%K t05-01-17 p
\textsc{Іван}. --- Аякже, запам'ятаю, бо попрошу в дядька маленьку кельму і
ночви, а висок зроблю сам із мотузочки.

%%K t05-tit-3 sub
\VUsaut{3.0mm}
\VUpk{10.2pt}{40}{Матеріали.}
\VUsaut{3.0mm}

%%K t05-01-18 p
\textsc{Дядько}. --- От і добре. Але скажи-но, Іване, що потрібно, щоб збудувати
дім?

%%K t05-01-19 p
\textsc{Іван}. --- Потрібен \VUgras{тесаний камінь} \textsuperscript{(59)} і
\VUgras{розчин} \textsuperscript{(65)}.

%%K t05-01-20 p
\textsc{Дядько}. --- Саме так. Поглянь на того чоловіка у великому капелюсі, на
його \VUgras{возі} \textsuperscript{(136)}. Це \VUgras{ломовий візник}
\textsuperscript{(135)}. Щодня він привозить сюди \VUgras{брили каменю}
\textsuperscript{(60)}. Він розвантажує воза на подвір'ї, а \VUgras{каменярі}
\textsuperscript{(83)} обтісують брили і пиляють їх своєю \VUgras{кам'яною
пилкою} \textsuperscript{(85)}. \VUgras{Підсобний робітник}
\textsuperscript{(146)} носить їх мулярові, а той кладе їх на місце.

%%K t05-01-21 p
Тим часом \VUgras{розчинник} \textsuperscript{(87)} готує розчин. Йому потрібні
\VUgras{пісок} \textsuperscript{(62)}, \VUgras{вапно} \textsuperscript{(63)} і
\VUgras{вода} \textsuperscript{(64)}. Вапно стоїть поруч із ним у \VUgras{мішку}
\textsuperscript{(71)}; \VUgras{мулярів підручний} \textsuperscript{(111)} везе
йому пісок у \VUgras{тачці} \textsuperscript{(112)}, а \VUgras{учень}
\textsuperscript{(113)} стоїть коло помпи \textsuperscript{(58)}. Він наповнює
\VUgras{відро} \textsuperscript{(114)}. Розчин скоро буде готовий.
\VUgras{Подавальник} \textsuperscript{(107)} бере його і несе сходами на
риштування, де муляр докінчує цей бік дому.

%%K t05-01-22 p
А тепер: як зветься робітник, який робить \VUgras{дах} \textsuperscript{(31)}?

%%K t05-01-23 p
\textsc{Іван}. --- Це \VUgras{покрівельник} \textsuperscript{(118)}.

%%K t05-tit-4 sub
\VUsaut{3.0mm}
\VUpk{10.2pt}{40}{Дах дому.}
\VUsaut{3.0mm}

%%K t05-01-24 p
\textsc{Пан Блан}. --- Так, але спершу йде \VUgras{тесля}
\textsuperscript{(115)}.

%%K t05-01-25 p
Тесля робить \VUgras{крокви} \textsuperscript{(35)}. Він з'єднує \VUgras{балки}
\textsuperscript{(37)} \VUgras{тиблями} \textsuperscript{(117)}. Ті робітники
вгорі, в широких оксамитових штанях, --- теслі. Один тримає брусок, другий теше
тибель своєю \VUgras{сокирою} \textsuperscript{(116)}. Той, що коло
\VUgras{димаря} \textsuperscript{(41)}, --- покрівельник. Він прибиває лати до
\VUgras{балок} (*), а \VUgras{шифер} \textsuperscript{(66)} --- до лат. Інколи
він кладе черепицю.

%%K t05-noto-1 noto
\VUnotes{143.2mm}{%
(*) \VUgras{Balk-o} = нім. \textit{Balken}, англ. \textit{joist}, фр.
\textit{solive}, іт. \textit{travicello}, укр. \textit{балка}, ісп. і порт.
\textit{viga}. Технічний корінь, ще не прийнятий, але пропонований тут для
передачі змісту французького слова \textit{solive}, яке не є ні
\textit{trabo} (фр. \textit{poutre}), ні брусок. Це обтесана колода, що несе
підлогу і стелю.%
}

%%K t05-01-26 p
Дах стайні \VUgras{черепичний} \textsuperscript{(55)}, дах дому
\VUgras{шиферний} \textsuperscript{(32)}, а дах веранди \VUgras{цинковий}
\textsuperscript{(24)}.

%%K t05-01-27 p
\textsc{Іван}. --- Цинкові дахи теж робить покрівельник?

%%K t05-01-28 p
\textsc{Пан Блан}. --- Ні, це \VUgras{бляхар} \textsuperscript{(121)} або
\VUgras{водопровідник}, он той, що вставляє \VUgras{слухове вікно}
\textsuperscript{(38)} в дах дому. Він же кладе \VUgras{ринви}
\textsuperscript{(43)} і \VUgras{водостічні труби} \textsuperscript{(42)}. Він
паяє цинк і бляху. Це він закріпив \VUgras{громовідвід} \textsuperscript{(40)} і
\VUgras{флюгер} \textsuperscript{(39)} на \VUgras{фронтоні}
\textsuperscript{(36)}.

%%K t05-01-29 p
\textsc{Іван}. --- Отже, дім готовий?

%%K t05-tit-5 sub
\VUsaut{3.0mm}
\VUpk{10.2pt}{40}{Інструменти.}
\VUsaut{3.0mm}

%%K t05-01-30 p
\textsc{Пан Блан}. --- Не зовсім. \VUgras{Тинькар} \textsuperscript{(99)} кладе
з \VUgras{цегли} \textsuperscript{(61)} перегородки, які ділять його на різні
кімнати. Він кладе \VUgras{тиньк} \textsuperscript{(70)} на \VUgras{стелі}
\textsuperscript{(26)} і на стіни. Зараз він працює над стелею \VUgras{веранди}
\textsuperscript{(33)}. Як і в муляра, у нього є \VUgras{кельма}
\textsuperscript{(100)} і \VUgras{ночви} \textsuperscript{(101)}.

%%K t05-01-31 p
Тоді столяр робить двері, вікна, \VUgras{дерев'яні підлоги}
\textsuperscript{(16)} і \VUgras{віконниці} \textsuperscript{(19)}.

%%K t05-01-32 p
Зараз він працює на подвір'ї за своїм \VUgras{верстатом} \textsuperscript{(90)}.
У нього є \VUgras{пилка} \textsuperscript{(94)}, щоб різати \VUgras{дерево}
\textsuperscript{(69)}, \VUgras{рубанок} \textsuperscript{(91)},
\VUgras{струбцина} \textsuperscript{(92)}, \VUgras{долото} \textsuperscript{(94
\textit{bis})}, \VUgras{циркуль} \textsuperscript{(95)} і дерев'яна
\VUgras{киянка} \textsuperscript{(95 \textit{bis})}.

%%K t05-01-33 p
\textsc{Іван}. --- Ах, та столяром бути ще веселіше, ніж муляром! Дядьку,
купиш мені верстат, пилку і рубанок? Я зроблю буду для Сірка.

%%K t05-01-34 p
\textsc{Дядько}. --- Куплю, мій хлопче.

%%K t05-01-35 p
\textsc{Іван}. --- Який я радий! А хто той чоловік, що стоїть коло вхідних
дверей?

%%K t05-tit-6 sub
\VUsaut{3.0mm}
\VUpk{10.2pt}{40}{Малярні роботи.}
\VUsaut{3.0mm}

%%K t05-01-36 p
\textsc{Пан Блан}. --- Це \VUgras{маляр} \textsuperscript{(102)}. Він стоїть на
своїй драбині з горщиком \VUgras{фарби} \textsuperscript{(103)} і своїм
\VUgras{пензлем} \textsuperscript{(104)}. Він укриває дерево вхідних дверей, щоб
воно довше трималося.

%%K t05-01-37 p
Він же обклеює кімнати шпалерами.

%%K t05-01-38 p
\VUgras{Шпалерник} \textsuperscript{(107)} на \VUgras{драбинці}
\textsuperscript{(106)}, на \VUgras{балконі} \textsuperscript{(28)}, вішає
\VUgras{штору} \textsuperscript{(29)}.

%%K t05-01-39 p
Робітник, що стоїть коло великого вікна, --- \VUgras{скляр}
\textsuperscript{(96)}. Він вставляє \VUgras{шибки} \textsuperscript{(18)} на
\VUgras{замазці} \textsuperscript{(73)}. А той робітник, що стоїть навколішки
коло підвальних дверей, --- \VUgras{слюсар} \textsuperscript{(97)}. Він ставить
\VUgras{замок} \textsuperscript{(6)}. \VUgras{Терпуг} \textsuperscript{(98)}
лежить поруч із ним.

%%K t05-01-40 p
\textsc{Іван}. --- А той, що б'є \VUgras{молотом} \textsuperscript{(84)} по
шматку заліза, теж слюсар?

%%K t05-01-41 p
\textsc{Пан Блан}. --- Точніше, він \VUgras{коваль} \textsuperscript{(125)}. Він
кує \VUgras{залізні частини} \textsuperscript{(67)} для \VUgras{електрика}
\textsuperscript{(124)}, який стоїть на даху \VUgras{возівні}
\textsuperscript{(51)}. У нього є \VUgras{горно} \textsuperscript{(126)},
\VUgras{ковадло} \textsuperscript{(127)} і \VUgras{обценьки}
\textsuperscript{(128)}.

%%K t05-01-42 p
Електрик закріплює \VUgras{мідний дріт} \textsuperscript{(44)} телефона.

%%K t05-tit-7 sub
\VUpk{10.2pt}{40}{Робота в саду.}
\VUsaut{3.0mm}

%%K t05-01-43 p
\textsc{Дядько}. --- У глибині \VUgras{двору} \textsuperscript{(45)}, коло
\VUgras{ґрат} \textsuperscript{(47)} і \VUgras{воріт} \textsuperscript{(48)}, ти
бачиш \VUgras{садівника} \textsuperscript{(129)}. Він готує маленький
\VUgras{сад} \textsuperscript{(49)}. Він скопав землю своїм \VUgras{заступом}
\textsuperscript{(131)}, сіє насіння і розрівнює \VUgras{граблями}
\textsuperscript{(130)}. Потім зі своєї \VUgras{поливальниці}
\textsuperscript{(132)} він поллє \VUgras{грядки} \textsuperscript{(150)}, і
рослини підуть у ріст.

%%K t05-01-44 p
\VUgras{Конюх} \textsuperscript{(133)}, який щойно прибрав коня і задав йому
корму, чистить \VUgras{екіпаж} \textsuperscript{(52)} губкою, \VUgras{ручною
помпою} \textsuperscript{(134)} і відром води. Потім він прибере його до
возівні і замкне двері важким \VUgras{засувом} \textsuperscript{(53)}.

%%K t05-01-45 p
Ну от, тепер ти, сподіваюся, задоволений; пояснень було досить.

%%K t05-01-46 p
\textsc{Іван}. --- Так, дядьку, але мені хотілося б подивитися дім усередині.

%%K t05-01-47 p
\textsc{Дядько}. --- Це буде завтра. Пан Ру пояснить тобі план і назве кожну
кімнату.

%%K t05-tit-8 sub
\VUsaut{3.0mm}
\VUpk{10.2pt}{40}{Розташування дому.}
\VUsaut{2.4mm}

%%K t05-apar-2 apar
{\centering\textit{(Див. план.)}\par}
\VUsaut{2.4mm}

%%K t05-01-48 p
\textsc{Архітектор}. --- Ходімо, Іване, я проведу тебе різними частинами дому.

%%K t05-01-49 p
Увійдімо в \VUgras{нижній поверх} \textsuperscript{(7)}. Ось кнопка
\VUgras{дзвінка} \textsuperscript{(11)}. Підіймаємося на три \VUgras{східці}
\textsuperscript{(14)} \VUgras{ґанку} \textsuperscript{(9)}, переступаємо
\VUgras{поріг} \textsuperscript{(10)} \VUgras{вхідних дверей}
\textsuperscript{(8)} --- і ось ми біля підніжжя \VUgras{сходів}
\textsuperscript{(13)}.

%%K t05-01-50 p
\textsc{Іван}. --- Це коридор.

%%K t05-01-51 p
\textsc{Архітектор}. --- Ні, це \VUgras{передпокій} \textsuperscript{(12)}.
\VUgras{Коридор} \textsuperscript{(a)} у глибині. Праворуч ---
\VUgras{вітальня} \textsuperscript{(b)} і \VUgras{їдальня}
\textsuperscript{(c)}; навпроти їдальні --- \VUgras{кухня}
\textsuperscript{(d)} і \VUgras{буфетна} \textsuperscript{(e)}; тоді, з боку
фасаду, \VUgras{кабінет} \textsuperscript{(f)}. \VUgras{Веранда}
\textsuperscript{(23)} зі своїми \VUgras{заскленими дверима}
\textsuperscript{(25)} прилягає до вітальні.

%%K t05-01-52 p
Тепер підіймімося сходами. Тримайся за \VUgras{поручні} \textsuperscript{(15)} і
лічи східці. Їх чотирнадцять. Ось \VUgras{площадка} \textsuperscript{(m)}: ми на
першому \VUgras{поверсі} \textsuperscript{(27)}.

%%K t05-tit-9 sub
\VUsaut{3.0mm}
\VUpk{10.2pt}{40}{Перший поверх.}
\VUsaut{3.0mm}

%%K t05-01-53 p
Навпроти сходів --- \VUgras{вбиральня} \textsuperscript{(20)} і \VUgras{туалетна
кімната} \textsuperscript{(j)}. Праворуч прекрасна \VUgras{спальня}
\textsuperscript{(g)} з двома вікнами на \VUgras{фасад} \textsuperscript{(1)}
дому. Ліворуч --- \VUgras{ванна} \textsuperscript{(1)} і \VUgras{кімната для
гостей} \textsuperscript{(i)} з великою \VUgras{альковою}
\textsuperscript{(h)}. Поруч зі спальнею --- \VUgras{дитяча}
\textsuperscript{(k)}. Вона виходить на широку \VUgras{терасу}
\textsuperscript{(21)}. Під цією терасою інша вбиральня \textsuperscript{(20)}, а
коло стіни --- \VUgras{дзвін} \textsuperscript{(22)}, у який дзвонять на обід.

%%K t05-01-54 p
\textsc{Іван}. --- Підіймімося на \VUgras{другий поверх}.

%%K t05-01-55 p
\textsc{Архітектор}. --- Другого поверху немає. Під дахом тільки
\VUgras{мансарди} \textsuperscript{(33)} і горище \textsuperscript{(34)}. Обхід
скінчено, можна спускатися.

%%K t05-01-56 p
\textsc{Іван}. --- Дякую, пане, було дуже цікаво. Я розкажу батькові все, що
бачив.
\VUsaut{4.0mm}
\VUfilet{22mm}
